\section{Introdução}

A crescente dependência das organizações em relação aos ativos de Tecnologia da Informação e Comunicação (TIC) transformou a TI de uma função meramente operacional em um componente estratégico central para a sobrevivência e competitividade empresarial \cite{weill2004governance}. Nesse cenário, a governança corporativa de TI estabelece os mecanismos de liderança, estruturas organizacionais e processos que garantem que as decisões e investimentos em tecnologia sustentem e ampliem os objetivos da organização.

O framework \textit{COBIT 2019} (\textit{Control Objectives for Information and Related Technologies}), desenvolvido pela ISACA, consolidou-se internacionalmente como uma das referências mais abrangentes para a governança e gestão de informações e tecnologia \cite{isaca2019framework}. Diferente de modelos puramente prescritivos, o COBIT 2019 destaca-se pela introdução de fatores de design (\textit{Design Factors}) e princípios customizáveis, permitindo desenhar sistemas de governança sob medida para as necessidades e riscos específicos de cada negócio \cite{isaca2019design}.

\subsection{Objetivos}

\subsubsection{Objetivo Geral}

Analisar a aplicação dos princípios e objetivos de governança e gestão do COBIT 2019 na estruturação de um modelo de Governança de TI corporativo.

\subsubsection{Objetivos Específicos}

\begin{itemize}
    \item Compreender a distinção conceitual e prática entre governança e gestão segundo o COBIT 2019;
    \item Explorar os 5 domínios fundamentais (EDM, APO, BAI, DSS e MEA) e seus respectivos objetivos;
    \item Analisar o papel dos fatores de design na personalização do sistema de governança;
    \item Avaliar a medição de capacidade e maturidade dos processos de governança.
\end{itemize}
